% ============================================================
%  Niels Treschow, "Om Gud, Idee- og Sandseverdenen samt de førstes
%    Aabenbarelse i den sidste. Et philosophisk Testament." Første Bog.
%    Christiania: Chr. Grøndahl, 1831.
%  TRANSLATION (English) — Hans Halvorson
%  SELECTION: vol. I, printed pp. 202–207 and 220–221.
%
%  Source of truth: transcription.tex. Mirrors it 1:1 (same 8 printed-page
%  markers, same two span headings, same emphasis spans).
%
%  TERMINOLOGY — continuous with ../enslige-ting/ and ../almindelig-logik/:
%    enslig / Ensligvæsen        -> singular / singular being
%    Individ / Individualitet    -> individual / individuality
%    besynderlig                 -> particular (18c sense)
%    Slægt / Art / Classe / Orden-> genus / species / class / order
%    Idee                        -> Idea (capitalised: Treschow's eternal Ideas)
%    Grundform                   -> fundamental form
%    Sandseverden                -> the sense-world
%    Vexelvirkning               -> reciprocity (reciprocal action)
%    Heelt / Enhed               -> whole / unity
%    Udflydelse                  -> emanation (in the neutral sense: outflow)
%    Forestilling                -> representation
%    Forstand                    -> understanding
%    Kraftyttring                -> manifestation of force
%    Charactermaler              -> portrayer of character
% ============================================================
\documentclass[12pt,a4paper]{article}
\usepackage[utf8]{inputenc}
\usepackage[T1]{fontenc}
\usepackage{libertinus}
\usepackage{libertinust1math}
\usepackage[english]{babel}
\usepackage[protrusion=true,expansion=true]{microtype}
\usepackage{setspace}
\usepackage[margin=1.2in]{geometry}
\usepackage{fancyhdr}
\usepackage{hyperref}

\hypersetup{
  pdftitle={On God, the World of Ideas and the World of Sense (1831) — extracts},
  pdfauthor={Niels Treschow},
  hidelinks
}

\pagestyle{fancy}
\fancyhf{}
\rhead{Treschow, \emph{On God} I (1831), extracts}
\cfoot{\thepage}

\setstretch{1.3}

\newcommand{\spanhead}[1]{%
  \bigskip\begin{center}\rule{0.18\linewidth}{0.4pt}\\[6pt]
  {\small\itshape #1}\\[4pt]\rule{0.18\linewidth}{0.4pt}\end{center}\smallskip}

\title{\textbf{On God, the World of Ideas and the World of Sense}\\[4pt]
  \large and the Revelation of the First in the Last\\[6pt]
  \normalsize A Philosophical Testament. First Book.\\[8pt]
  \large Extracts: printed pp.\ 202--207 and 220--221}
\author{Niels Treschow\\[4pt]
  \small Translated from the Danish by Hans Halvorson}
\date{\small Christiania: Chr.\ Grøndahl, 1831}

\begin{document}
\maketitle
\thispagestyle{fancy}

\bigskip

\spanhead{Extract A: printed pp.\ 202--207.\\
  The individuality of the Ideas; the indeterminacy of genus- and
  species-concepts.}

% printed p. 202
\noindent
[\dots] who, without compass or visible stars, are tossed about upon the sea, and
therefore gladly make land at the nearest shore to which fortune carries them.
Should there then not be some principle that can determine us on this voyage of
inquiry, and at the same time, like a fixed pole, guide our course? Is it
improbable---and not rather a conception which the nature of the case must itself
bring with it---that, just as the highest being is necessarily active in the
production of the countless forms or Ideas in which it reveals itself, so these
must by a like necessity do the same? Now they, as actual things, cannot be
anything other than individual; for only individuals, not genera and species, are
in every respect so determinate that they can exist for themselves. If, then, the
Ideas---as one must suppose from the perfection of their nature and their close
kinship with the divine---can also create in imitation of the latter, then this
creation must likewise be something individual: namely a form of themselves again,
one that continually develops, not several different forms. So far their power
cannot extend, since it, though in one respect infinite, is yet itself only a
singular form. In such an individual the fundamental form---that is, its
essence---alone is abiding and unchangeable; in the accidental forms which it
gradually assumes in the course of its advance there can be no permanence, for
these are possible only in time, where
% printed p. 203
sheer alternation reigns. Thus a sense-world arises out of the ideal one, and with
the same necessity as that with which the latter has itself sprung from an
absolute unity and trinity.

Every Idea accordingly brings forth its own image, but in countless momentarily
changing shapes, from the germ up to the ripened fruit. Nor is all further
development over with this. Death itself, which follows upon the fruit, is itself
only a transformation. In perennial plants every spring drives forth fresh shoots.
The tree grows as long as it lives, and when at length it dies, it continues in its
offspring an existence altered in many respects, perhaps in the end even ennobled.
The same is to be supposed of soft herbs and annual growths. The Idea, of which
everything sensible is an expression, is itself infinite, but the image is finite,
so that the latter, however much it hastens, however diligently it works, never
reaches the goal of its striving. All of us, the inhabitants of this sense-world,
stand to our Ideas as these do to the absolutely One. Both kinds of being are
outflowings of the same, but in a different way: partly immediate, partly mediate.
So too does the human understanding lay its plans; its ideas, which may be of
countless kinds---of well-cultivated fields, pleasure gardens, magnificent
buildings, scientific theories, works of art, public and private arrangements for
the use of life---are each
% printed p. 204
for itself present all at once in the mind of their author; but outwardly they are
only very slowly brought anywhere near that perfection which one beholds at first
only obscurely, and then, as the work proceeds, ever more clearly.

If one would say that it is here the understanding, not the Idea, which brings
forth or creates, then the understanding---although one commonly seems to wish to
present it as an actual force---is nevertheless nothing other than an abstract and
general concept, under which countless manifestations of force, or forces properly
so called, are comprehended. These now are the human Ideas, and the many sorts of
higher representations that arise from them: the same are consequently also in
truth the only productive ones.

In all this it must be noted that I am obliged here to call the higher
representations---which are regarded as offspring of reason, but which, properly
speaking, develop out of a single one, as out of a noble shoot, like those just
adduced---by the same name as those eternal forms of the absolutely One. The analogy
between God and the human spirit, the uncreated and the created, the infinite and
the finite, warrants this manner of speaking.

I said just now that no eternal Idea brings forth more than its own individual
image. Against this it seems to tell that the human understanding in the same
person or individual can nevertheless bring forth many quite different plans and
works, which cannot be regarded as mere forms or images of him
% printed p. 205
self. Thus a work of art, and indeed even the plan laid for it in thought, is no
image of the artist's spirit; and even were one of the many to be so regarded, it
still could not hold of the rest, which often have very little resemblance to him
or to one another. Nonetheless every \emph{portrayer of character} knows---whether
he would depict some person with the brush, in wax, in marble, or in words and
speech---that his whole appearance, all his ideas, plans and actions are the
singular traits which together constitute his individuality; and if any of these
is lacking, it is not presented in its wholeness as it ought to be. Nor must the
diversity in the content of the works, or of the objects produced, be taken into
account here, but rather a certain unmistakable likeness in the form, whereby
everything is marked out that is to be ascribed to the same individual. Further,
individuality is composed of many parts, inward as well as outward, that may be
noticed. It must in every case be as difficult, or more so, to derive these from
any unity of the fundamental form, as it is in human nature generally to find out
the fundamentally essential on which all the likewise abiding properties that
express themselves in it depend. Nonetheless there is no doubt that every singular
being---however great a manifoldness one may discover in him, as in his eternal
Idea---is yet only one. The objection that from one Idea there may proceed countless
ones, which are as little
% printed p. 206
mere forms of the same individual as human plans and works are of the individual
understanding, may accordingly be answered thus: that the latter is nevertheless
really so, and that in the different individuals in whose production this
understanding has some share, only certain traits belong to it, which taken
together constitute not several, but only a single image of itself. Every such
trait is a particular form, but not one subsisting by itself; it subsists solely
in and through the objects in which it is expressed. Between the human
understanding and the eternal Ideas there is accordingly this significant
difference: that the former can bring forth nothing subsisting for itself. What
gives its effects a certain permanence is not present by means of it, but was there
along with it---namely what one usually calls the matter of things. The Ideas, by
contrast, can, like God, and by means of an activity granted them by him for that
purpose, bring forth everything that in the sense-world one takes for
self-subsistent, or for an actual substance.

The theory here advanced differs essentially from the Platonic in this, that the
latter makes the Ideas into the objects of general concepts, or conversely makes
these into Ideas---among which not only those genera and species must be reckoned
under which all natural singular beings, such as singular crystals, plants and
animals, are comprehended, but even mere products of art, such as tables, chairs,
and so on. \emph{Plato}, that is, regarded all
% printed p. 207
individuals as things in which there is nothing wholly abiding, and knew of no
unchangeable fundamental form, gradually developing, to be spoken of in them---
although it is clear that \emph{Socrates} and \emph{Alcibiades} just as much remain
for all eternity Socrates and Alcibiades as the animal an animal, the plant a
plant; indeed the former---as I shall explain more closely in a separate division
of the following book---so much the more, since all \emph{class-, genus- and
species-concepts are more or less wavering and indeterminable, whereas the
individuals are in themselves so perfectly determinate that a change or transition
from one to another is impossible}, while an individual, by contrast, may very well
pass over into another species, genus, order or class than the one to which it at
present belongs. Either there are no Ideas at all---though their existence has been
proved---or these consist in individual, eternal, unchangeable and perfect
fundamental forms, which are expressed from the ideal world in the sense-world and
are there ever further developed. Now if the Ideas are individual, then there can
be no other difference among them than in respect of this their form; but in
perfection all must, each after its own kind, be equal. They do not, then, stand
at different stages of development, nor are they lower and higher in the same sense
in which genera, species and individuals are said to stand in such a relation to
one another. There must, on the contrary, among the Ideas, as in the sense-world,
be both indi\-[viduals \dots]

\spanhead{Extract B: printed pp.\ 220--221.\\
  The unity of the ideal world by reciprocity; individuals as the constituents
  of the sense-world.}

% printed p. 220
\noindent
[\dots] In the ideal world, as in the real, a constant reciprocity must obtain:
otherwise it could constitute no actual unity, no whole. Its members must be able
to communicate themselves, to know one another, without their state being thereby
altered; for they do not live in time, in which all change takes place. This
communication is more or less mediate according to the natural relation in which
they stand to one another as equal and unequal, superior and subordinate. This
relation is abiding; and from it there arises at the same time a difference of
activity with respect to nearer and remoter objects. So far these properties of the
Ideas differ from the divine omnipotence
% printed p. 221
and omniscience, in whose infinity mediacy and distance vanish altogether. A
precisely similar relation may be observed among the human ideas of reason: they
are all of them more and less akin to one another; of whatever kind---moral,
political, physical, economic---they are not without reciprocal influence, and
they constitute, once duly ordered, a coherent whole. From ourselves, from our own
consciousness, as from the only source accessible to us, we must draw everything
that we desire and can hope to learn about the higher world. According to the
principles of the identity-system we shall not be deceived in this either; for
essentially the highest and the lowest things are entirely alike. But besides the
activity common to them, every Idea possesses also its own peculiar one, which
consists in the production and development of an image answering to its individual
essence under the ever-changing forms of time and space. These are the individuals
of which the sense-world consists. The Ideas form for themselves in this way a kind
of body, by means of which---yet without being affected and altered thereby, as by
something outward---they put themselves in a position to behold their own inward
being with full consciousness from every possible side; for therein they can pursue
it through all the stages of a possible existence, and not merely as it is at the
very highest, at which they actually stand. In every sensible individual its Idea
thus reveals itself as in an image ever proper to it [\dots]

\begin{center}\rule{0.25\linewidth}{0.4pt}\end{center}

\end{document}
